% -----------------------------------------------------------
\subsection{Flujos de Negocio}
% -----------------------------------------------------------

Esta sección describe los principales flujos operativos del sistema Bouquet. Los
flujos se organizan de acuerdo con los actores participantes y con las entidades
centrales del modelo: usuarios colaboradores, mesas físicas, sesiones virtuales,
órdenes, cocina, liquidación interna y procesamiento analítico.

\subsubsection{Login de Usuario Colaborador}

Roles participantes: Admin Cadena, Gerente Zona, Gerente Sucursal, Mesero y
Cocina/Chef.

El usuario colaborador ingresa su correo electrónico y contraseña en la pantalla
de inicio de sesión. El sistema valida las credenciales y consulta los roles
contextuales asociados mediante \texttt{UserRole}. Con base en el rol y alcance
organizacional del usuario, la aplicación redirige a la interfaz correspondiente.

En caso de primer acceso, el sistema puede solicitar actualización de contraseña
antes de permitir el uso normal de la plataforma. Una vez autenticado, el usuario
solo podrá consultar o modificar información dentro del alcance permitido por su
rol: cadena, zona o sucursal.

\subsubsection{Alta de Zonas}

Roles participantes: Admin Cadena.

El Admin Cadena accede al módulo de zonas desde el panel administrativo. El sistema
muestra las zonas existentes de la cadena y permite registrar nuevas zonas operativas.
Cada zona queda asociada a una cadena restaurantera y sirve como agrupación geográfica
o administrativa de sucursales.

\subsubsection{Alta de Sucursales}

Roles participantes: Admin Cadena y Gerente Zona.

El Admin Cadena puede crear sucursales dentro de cualquier zona perteneciente a la
cadena. El Gerente Zona únicamente puede crear o administrar sucursales dentro de
las zonas que tiene asignadas. Cada sucursal define su moneda de operación, parámetros
locales, configuración de pedidos móviles y reglas básicas de operación.

\subsubsection{Alta de Categorías, Productos y Modificadores}

Roles participantes: Admin Cadena, Gerente Zona y Gerente Sucursal.

Los usuarios autorizados pueden administrar el menú de acuerdo con su alcance
organizacional. Las categorías agrupan productos; los productos pueden tener variantes,
modificadores, precio local, grupo fiscal y estación de preparación asignada.

Cuando un producto, variante o modificador cambia, las órdenes históricas no se alteran,
debido a que \texttt{OrderItem} y \texttt{OrderItemModifier} conservan snapshots de
nombre, precio, tasa de impuesto, estación y ajustes aplicados al momento de ordenar.

\subsubsection{Alta de Usuarios con Roles}

Roles participantes: Admin Cadena, Gerente Zona y Gerente Sucursal.

El usuario creador registra el correo del nuevo colaborador, sus datos básicos y el rol
correspondiente. El sistema valida que el creador tenga permisos suficientes y que el
rol asignado pertenezca a su alcance operativo.

El Admin Cadena puede crear usuarios dentro de la cadena. El Gerente Zona puede crear
usuarios dentro de sus sucursales asignadas. El Gerente Sucursal puede crear personal
operativo de su sucursal, como meseros y usuarios de cocina. Toda creación de usuario
queda registrada en \texttt{AuditLog}.

\subsubsection{Creación y Gestión de Mesa Física}

Roles participantes: Gerente Sucursal, Mesero.

Desde el módulo de mesas, el usuario autorizado puede crear, editar o liberar mesas
físicas. Cada \texttt{DiningTable} pertenece a una sucursal y contiene número, capacidad,
código público para QR, versión de QR y estado físico-operativo.

Los estados posibles de una mesa física son: \texttt{LIBRE}, \texttt{OCUPADA} y
\texttt{NO\_DISPONIBLE}. Una mesa en estado \texttt{LIBRE} puede asociarse a una nueva
sesión virtual; una mesa \texttt{OCUPADA} no puede formar parte de otra sesión activa.

\subsubsection{Inicialización de Mesa Virtual}

Roles participantes: Mesero, Gerente Sucursal.

Desde la pantalla de mesas, el usuario selecciona una o varias mesas físicas en estado
\texttt{LIBRE}. El sistema valida que las mesas pertenezcan a la sucursal del usuario y
que no estén asociadas a otra \texttt{DiningSession} activa.

Si la validación es exitosa, el sistema crea una nueva \texttt{DiningSession} en estado
\texttt{PENDIENTE\_ACTIVACION}, registra las asociaciones en \texttt{DiningSessionTable}
y cambia las mesas físicas a estado \texttt{OCUPADA}. Posteriormente genera el enlace QR
firmado con HMAC-SHA256 a partir de \texttt{DiningTable.publicCode} y
\texttt{DiningTable.qrVersion}.

\subsubsection{Acceso de Comensales a la Mesa Virtual}

Roles participantes: Anfitrión, Comensal.

El flujo de acceso se divide en dos subflujos:

\begin{enumerate}[leftmargin=1.5em]
  \item \textbf{Activar mesa virtual:} el Anfitrión escanea el QR, valida el enlace
  firmado y crea una contraseña de acceso para la mesa. El sistema almacena únicamente
  el hash de la contraseña en \texttt{DiningSession.accessCodeHash}.

  \item \textbf{Unirse a mesa virtual:} los demás comensales ingresan mediante el código
  de sesión y la contraseña compartida por el Anfitrión. Al unirse, el sistema crea un
  registro \texttt{Guest} asociado a la \texttt{DiningSession}.
\end{enumerate}

El sistema no almacena la contraseña en texto plano ni la muestra nuevamente después
de su creación. El acceso de comensales es temporal y se limita al ciclo de vida de la
sesión virtual.

\subsubsection{Diagrama de Secuencia --- Flujo de Activación de Mesa Virtual}

La Figura~\ref{fig:seq_activacion} muestra el flujo de activación de una mesa virtual
por parte del Anfitrión. Este flujo es crítico porque valida el QR firmado, crea la
contraseña de acceso y cambia el estado de la sesión de \texttt{PENDIENTE\_ACTIVACION}
a \texttt{ACTIVA}.

\begin{figure}[ht]
\centering
\small
\fbox{
\begin{minipage}{0.92\textwidth}
\centering
\textbf{Diagrama de Secuencia --- Activación de Mesa Virtual}\\[8pt]
\begin{tabular}{>{\raggedright\arraybackslash}p{13cm}}
\textbf{Actores:} :Mesero, :Next.js API, :Prisma, :PostgreSQL, :Supabase Realtime, :Anfitrión \\
\end{tabular}\\[4pt]
\begin{tabular}{>{\raggedright\arraybackslash}p{13cm}}
\textbf{1.} Mesero $\rightarrow$ API: \texttt{POST /api/dining-sessions} \{tableIds[]\} \\
\textbf{2.} API: validar sesión del mesero y permisos de sucursal \\
\textbf{3.} API $\rightarrow$ PostgreSQL: verificar mesas físicas en estado \texttt{LIBRE} \\
\textbf{4.} PostgreSQL $\rightarrow$ API: mesas disponibles \\
\textbf{5.} API $\rightarrow$ PostgreSQL: crear \texttt{DiningSession(status=PENDIENTE\_ACTIVACION)} \\
\textbf{6.} API $\rightarrow$ PostgreSQL: crear asociaciones en \texttt{DiningSessionTable} \\
\textbf{7.} API $\rightarrow$ PostgreSQL: actualizar \texttt{DiningTable.status = OCUPADA} \\
\textbf{8.} API: generar enlace QR firmado con HMAC-SHA256 usando \texttt{publicCode + qrVersion} \\
\textbf{9.} API $\rightarrow$ Mesero: QR generado para compartir con el Anfitrión \\[6pt]
\textit{--- El Anfitrión escanea el QR ---} \\[6pt]
\textbf{10.} Anfitrión $\rightarrow$ API: \texttt{POST /api/dining-sessions/activate} \{publicCode, signature, password\} \\
\textbf{11.} API: verificar firma HMAC, versión QR y estado de sesión \\
\quad\quad \textbf{[alt: firma inválida o mesa no asociada]} $\rightarrow$ HTTP~401/409 \\
\textbf{12.} API $\rightarrow$ PostgreSQL: transacción atómica \\
\quad\quad \texttt{hash(password)} \\
\quad\quad \texttt{actualizar DiningSession(status=ACTIVA, activatedAt=now())} \\
\quad\quad \texttt{crear Guest(isHost=true)} \\
\textbf{13.} PostgreSQL $\rightarrow$ API: sesión activada \\
\textbf{14.} API $\rightarrow$ Supabase Realtime: broadcast \texttt{dining\_session\_updated} \\
\textbf{15.} API $\rightarrow$ Anfitrión: HTTP~200 + \texttt{joinCode} + sesión temporal \\
\textbf{16.} Anfitrión: redirección al Home de mesa virtual \\
\end{tabular}
\end{minipage}
}
\caption{Diagrama de secuencia del flujo de activación de mesa virtual. La activación
valida un QR firmado con HMAC-SHA256 y actualiza la sesión a estado \texttt{ACTIVA}.}
\label{fig:seq_activacion}
\end{figure}

\subsubsection{Creación y Solicitud de Órdenes}

Roles participantes: Anfitrión, Comensal, Mesero.

El comensal navega el menú, selecciona productos, variantes y modificadores, y los agrega
a su canasta. Al confirmar, el sistema crea una \texttt{RestaurantOrder} y sus
\texttt{OrderItem}. Cada ítem conserva snapshots de producto, precio, estación, impuesto
y modificadores aplicados.

Las órdenes pueden registrarse como consumo individual o compartido. Cuando se confirma
la primera orden de la sesión, la \texttt{DiningSession} cambia de \texttt{ACTIVA} a
\texttt{EN\_CONSUMO}, y se registra \texttt{firstOrderAt}.

\subsubsection{Flujo de Cocina}

Roles participantes: Cocina/Chef, Mesero.

El personal de cocina visualiza los ítems pendientes en el tablero KDS. Cada ítem puede
cambiar de estado siguiendo el flujo:

\[
\texttt{PENDIENTE} \rightarrow \texttt{EN\_PREPARACION} \rightarrow \texttt{LISTA}
\rightarrow \texttt{ENTREGADA}
\]

El personal de cocina actualiza los estados hasta \texttt{LISTA}. El mesero confirma la
entrega al comensal mediante el cambio de \texttt{LISTA} a \texttt{ENTREGADA}. Cada
transición se registra en \texttt{OrderItemStatusEvent}, lo que permite reconstruir el
flujo de cocina y calcular tiempos de preparación en la capa analítica.

\subsubsection{Flujo de Liquidación Interna}

Roles participantes: Anfitrión, Comensal, Mesero, Gerente Sucursal.

El comensal accede a la sección de liquidación interna y visualiza el desglose de consumo
individual y compartido. El sistema muestra el total interno de cuenta, aportaciones
registradas y saldo pendiente.

El comensal puede registrar una aportación interna asociada a su consumo, a ítems
específicos o a la mesa compartida. El sistema crea un registro
\texttt{SettlementContribution} y, cuando corresponde, registros
\texttt{SettlementAllocation} para distribuir la aportación sobre ítems o comensales.

Para evitar sobreliquidación, el registro de aportaciones se realiza dentro de una
transacción ACID. El sistema bloquea la fila de \texttt{Settlement}, valida
\texttt{remainingCents} y rechaza cualquier aportación que exceda el saldo pendiente.

Cuando \texttt{Settlement.remainingCents = 0}, la liquidación cambia a
\texttt{LIQUIDADA} y la sesión actualiza \texttt{DiningSession.status} a
\texttt{LIQUIDADA}. El sistema genera una constancia interna de estado de cuenta y
notifica al mesero para realizar el cierre operativo.

\subsubsection{Cierre Operativo de Mesa}

Roles participantes: Mesero, Gerente Sucursal.

Una vez liquidada la sesión, el mesero o gerente puede cerrar operativamente la mesa.
El sistema actualiza la \texttt{DiningSession} a estado \texttt{CERRADA}, registra
\texttt{closedAt}, actualiza las asociaciones \texttt{DiningSessionTable.leftAt} y cambia
las mesas físicas asociadas a estado \texttt{LIBRE}, salvo que alguna haya sido marcada
como \texttt{NO\_DISPONIBLE}.

El cierre operativo queda registrado en \texttt{AuditLog}.

\subsubsection{Ejecución del Pipeline Analítico}

Roles participantes: Sistema, Gerente Sucursal, Gerente Zona, Admin Cadena.

El pipeline analítico se ejecuta de forma programada o bajo solicitud autorizada desde el
dashboard. Cada ejecución genera un registro \texttt{AnalyticsJobRun}, indicando tipo de
job, estado, ventana de datos, tiempos de inicio y fin, errores y metadatos técnicos.

El proceso extrae datos operativos hacia la capa Bronze, transforma y limpia información
en Silver mediante Apache Spark, y genera tablas Gold para consulta del dashboard. Las
métricas históricas del dashboard no se calculan directamente sobre las tablas OLTP, sino
sobre agregados precalculados como \texttt{AnalyticSalesDaily},
\texttt{AnalyticItemVelocity}, \texttt{AnalyticServiceTimes} y
\texttt{AnalyticDemandEstimate}.